\documentclass{article}
\usepackage{amsmath,amssymb,amsthm}
\usepackage{algorithm}
\usepackage{algpseudocode}
\usepackage{hyperref}
\usepackage{enumitem}
\setlist{noitemsep}
\newtheorem{theorem}{Theorem}
\newtheorem{lemma}{Lemma}
\newtheorem{assumption}{Assumption}
\newcommand{\norm}[1]{\left\lVert#1\right\rVert}
\newcommand{\dualnorm}[1]{\left\lVert#1\right\rVert_{*}}
\newcommand{\inner}[2]{\left\langle #1,#2\right\rangle}
\newcommand{\Breg}[2]{D_{h}(#1,#2)}
\begin{document}

\section*{Mirror Descent with Bregman Prox (MD‑B)}

\section*{Mathematical Formulation}
Let $f:\mathbb{R}^{d}\rightarrow\mathbb{R}$ be differentiable and let $h:\mathbb{R}^{d}\rightarrow\mathbb{R}$ be a \emph{prox‑function}, i.e. a continuously differentiable, $1$‑strongly convex function with respect to a norm $\|\cdot\|$.
The Bregman divergence generated by $h$ is
\[
\Breg{x}{y}=h(x)-h(y)-\inner{\nabla h(y)}{x-y}\, .
\]

Given a stepsize $\eta_{k}>0$, the MD‑B iteration is
\begin{equation}\label{eq:update}
x_{k+1}
\;=\;
\arg\min_{x\in\mathbb{R}^{d}}
\Big\{ \inner{\nabla f(x_{k})}{x}
      +\frac{1}{\eta_{k}}\Breg{x}{x_{k}}\Big\}.
\end{equation}
The optimality condition of \eqref{eq:update} reads
\begin{equation}\label{eq:optcond}
\inner{\nabla f(x_{k})+\tfrac{1}{\eta_{k}}\bigl(\nabla h(x_{k+1})-\nabla h(x_{k})\bigr)}{x-x_{k+1}}\;\ge\;0,
\qquad\forall x\in\mathbb{R}^{d}.
\end{equation}
Define the \emph{mirror gradient mapping}
\[
G_{\eta_{k}}(x_{k}) \;:=\; \frac{1}{\eta_{k}}\bigl(\nabla h(x_{k})-\nabla h(x_{k+1})\bigr).
\]

\section*{Pseudocode}
\begin{algorithm}[H]
\caption{MD‑B (Mirror Descent with Bregman Prox)}
\begin{algorithmic}[1]
\Require Initial point $x_{0}\in\mathbb{R}^{d}$, stepsize schedule $\{\eta_{k}\}_{k\ge0}>0$, prox‑function $h$.
\For{$k=0,1,2,\dots$}
    \State Compute stochastic (here deterministic) gradient $g_{k}\gets\nabla f(x_{k})$.
    \State $x_{k+1}\gets\displaystyle\arg\min_{x}\Big\{ \inner{g_{k}}{x}+\frac{1}{\eta_{k}}D_{h}(x,x_{k})\Big\}$.
\EndFor
\State \textbf{output} $x_{K}$ or the averaged point $\bar x_{K}:=\frac{1}{K}\sum_{k=0}^{K-1}x_{k}$.
\end{algorithmic}
\end{algorithm}

\section*{Dry Run Example}
Consider the one–dimensional problem
\[
f(w)=(w-3)^{2},\qquad 
\nabla f(w)=2(w-3).
\]
Choose the Euclidean prox‑function $h(w)=\tfrac12 w^{2}$, for which
$D_{h}(x,y)=\tfrac12 (x-y)^{2}$ and $\nabla h(w)=w$.
Let $w_{0}=0$ and a constant stepsize $\eta=0.1$.
The subproblem at iteration $k$ becomes
\[
w_{k+1}
=\arg\min_{w}\Big\{2(w_{k}-3)w+\frac{1}{0.1}\tfrac12 (w-w_{k})^{2}\Big\}
=\arg\min_{w}\Big\{-2(3-w_{k})w+5(w-w_{k})^{2}\Big\}.
\]

\begin{tabular}{c|c|c|c}
\textbf{Iter.}&$w_{k}$&$\nabla f(w_{k})$&$w_{k+1}= \displaystyle\frac{2\eta}{1+2\eta}(3)+\frac{1}{1+2\eta}w_{k}$\\\hline
0&0&$-6$&$w_{1}=0.6$\\
1&0.6&$-4.8$&$w_{2}=1.08$\\
2&1.08&$-3.84$&$w_{3}=1.464$\\
\end{tabular}

Objective values:
\[
f(w_{0})=9,\; f(w_{1})=5.76,\; f(w_{2})\approx3.69,\; f(w_{3})\approx2.37,
\]
showing monotone decrease.

\section*{Assumptions}
\begin{assumption}[Strong Convexity of the Prox‑function]\label{as:strong}
$h$ is $1$‑strongly convex w.r.t. $\|\cdot\|$:
\[
h(y)\ge h(x)+\inner{\nabla h(x)}{y-x}+\tfrac12\|y-x\|^{2},
\qquad\forall x,y.
\]
Consequently,
\[
\Breg{x}{y}\ge\tfrac12\|x-y\|^{2}.
\]
\end{assumption}

\begin{assumption}[Smoothness of $f$ w.r.t. the Bregman geometry]\label{as:smooth}
There exists $L>0$ such that for all $x,y$,
\[
f(y)\le f(x)+\inner{\nabla f(x)}{y-x}+\frac{L}{2}\|y-x\|^{2}.
\]
\end{assumption}

\begin{assumption}[Bounded Level Set]\label{as:bound}
The set $\{x\mid f(x)\le f(x_{0})\}$ is bounded, and we denote
\[
\Delta_{0}:=f(x_{0})-\inf_{x}f(x)<\infty .
\]
\end{assumption}

\section*{Main Convergence Theorem}
\begin{theorem}\label{thm:main}
Assume \ref{as:strong}–\ref{as:bound} and choose a constant stepsize
$\eta\in\big(0,\;1/L\big]$.
Define the mirror gradient mapping $G_{\eta}(x_{k})$ as above.
Then for any integer $T\ge1$,
\[
\frac{1}{T}\sum_{k=0}^{T-1}\|G_{\eta}(x_{k})\|_{*}^{2}
\;\le\;
\frac{2\Delta_{0}}{\eta\,(1-L\eta)\,T}.
\]
Consequently,
\[
\min_{0\le k<T}\|G_{\eta}(x_{k})\|_{*}^{2}
\;=\;O\!\left(\frac{1}{T}\right).
\]
\end{theorem}

\section*{Theoretical Properties}
\begin{enumerate}
\item \textbf{Stability.} The update \eqref{eq:update} is a proximal step in the dual space; the strong convexity of $h$ guarantees that the iterates remain in a bounded region whenever the level set of $f$ is bounded (Assumption~\ref{as:bound}).
\item \textbf{Hyperparameter Sensitivity.} The stepsize must satisfy $\eta\le 1/L$ (cf. Theorem~\ref{thm:main}) to keep the coefficient $1-L\eta$ positive. Larger $\eta$ yields faster per‑iteration decrease but may violate the condition, leading to possible divergence.
\item \textbf{Comparison to Euclidean Gradient Descent.} When $h(x)=\tfrac12\|x\|^{2}$, $\Breg{x}{y}=\tfrac12\|x-y\|^{2}$ and $G_{\eta}(x_{k})=\tfrac{1}{\eta}(x_{k}-x_{k+1})$; the theorem reduces to the standard $O(1/T)$ rate for non‑convex smooth gradient descent.
\item \textbf{Special Cases.}
\begin{itemize}
\item If $f$ is convex, the same proof yields the classic $O(1/T)$ bound on the primal gap $f(\bar x_{T})-f^{*}$.
\item If $h$ is the negative entropy and $\|\cdot\|$ is the $\ell_{1}$ norm, MD‑B becomes the exponentiated gradient method; the theorem still holds with the $\ell_{\infty}$ dual norm.
\end{itemize}
\end{enumerate}

\section*{Discussion}
The theorem shows that even without convexity of $f$, the mirror gradient mapping vanishes at a rate $1/T$, which is the standard benchmark for first‑order methods on smooth non‑convex problems. The Bregman geometry allows the algorithm to adapt to the underlying structure of the feasible set (e.g., simplex, probability simplex) while preserving the same worst‑case guarantee as Euclidean gradient descent.

\section*{Preliminaries (Definitions and Lemmas)}
\begin{lemma}[Three‑point identity]\label{lem:three}
For any $x,y,z$,
\[
\Breg{x}{z}-\Breg{y}{z}-\Breg{x}{y}
=\inner{\nabla h(z)-\nabla h(y)}{x-y}.
\]
\end{lemma}
\begin{proof}
Direct expansion of the definition of $D_{h}$ yields the identity. \hfill\qedsymbol
\end{proof}

\begin{lemma}[Norm bound from strong convexity]\label{lem:sc}
If $h$ satisfies Assumption~\ref{as:strong}, then
\[
\Breg{x}{y}\ge\tfrac12\|x-y\|^{2},\qquad
\| \nabla h(x)-\nabla h(y) \|_{*}\ge \|x-y\|.
\]
\end{lemma}
\begin{proof}
The first inequality is exactly the definition of $1$‑strong convexity.
The second follows by differentiating the strong convexity inequality and using
the definition of the dual norm. \hfill\qedsymbol
\end{proof}

\section*{Main Proof (Step‑by‑Step Derivation)}
We prove Theorem~\ref{thm:main}. Throughout the proof we fix a constant stepsize
$\eta\in(0,1/L]$ and write $G_{k}=G_{\eta}(x_{k})$.

\begin{enumerate}
\item[Step 1.] \textbf{Optimality condition.}  
From \eqref{eq:optcond} with $x=x_{k}$ we obtain
\[
\inner{\nabla f(x_{k})}{x_{k}-x_{k+1}}
\;\ge\;
\frac{1}{\eta}\inner{\nabla h(x_{k+1})-\nabla h(x_{k})}{x_{k}-x_{k+1}}
= \eta\|G_{k}\|_{*}^{2}.
\tag{1}
\]
(The last equality uses the definition of $G_{k}$.)

\item[Step 2.] \textbf{Smoothness descent.}  
Using $L$‑smoothness (Assumption~\ref{as:smooth}) with $y=x_{k+1}$,
\[
f(x_{k+1})\le f(x_{k})+\inner{\nabla f(x_{k})}{x_{k+1}-x_{k}}
+\frac{L}{2}\|x_{k+1}-x_{k}\|^{2}.
\tag{2}
\]
Replace $\inner{\nabla f(x_{k})}{x_{k+1}-x_{k}}=-\inner{\nabla f(x_{k})}{x_{k}-x_{k+1}}$ and apply (1):
\[
f(x_{k+1})
\le f(x_{k})-\eta\|G_{k}\|_{*}^{2}
+\frac{L}{2}\|x_{k+1}-x_{k}\|^{2}.
\tag{3}
\]

\item[Step 3.] \textbf{Relating the norm to the gradient mapping.}  
By Lemma~\ref{lem:sc},
\[
\|x_{k+1}-x_{k}\|
\;\le\;
\|\nabla h(x_{k})-\nabla h(x_{k+1})\|_{*}
= \eta\|G_{k}\|_{*}.
\tag{4}
\]
Insert (4) into (3):
\[
f(x_{k+1})
\le f(x_{k})-\eta\|G_{k}\|_{*}^{2}
+\frac{L}{2}\eta^{2}\|G_{k}\|_{*}^{2}
= f(x_{k})-\bigl(\eta-\tfrac{L}{2}\eta^{2}\bigr)\|G_{k}\|_{*}^{2}.
\tag{5}
\]

\item[Step 4.] \textbf{Descent coefficient positivity.}  
Since $\eta\le 1/L$, we have $1-L\eta\ge0$ and therefore
\[
\eta-\tfrac{L}{2}\eta^{2}
= \eta\Bigl(1-\tfrac{L}{2}\eta\Bigr)
\;\ge\;
\frac{\eta}{2}\bigl(1-L\eta\bigr).
\tag{6}
\]

\item[Step 5.] \textbf{Recursive inequality.}  
Combining (5) and (6),
\[
f(x_{k+1})
\le f(x_{k})-\frac{\eta}{2}\bigl(1-L\eta\bigr)\|G_{k}\|_{*}^{2}.
\tag{7}
\]

\item[Step 6.] \textbf{Summation and telescoping.}  
Sum (7) over $k=0,\dots,T-1$:
\[
f(x_{T})\le f(x_{0})
-\frac{\eta}{2}\bigl(1-L\eta\bigr)\sum_{k=0}^{T-1}\|G_{k}\|_{*}^{2}.
\tag{8}
\]
Rearrange and use $f(x_{T})\ge\inf f$:
\[
\sum_{k=0}^{T-1}\|G_{k}\|_{*}^{2}
\le
\frac{2\bigl(f(x_{0})-\inf f\bigr)}{\eta\bigl(1-L\eta\bigr)}
=
\frac{2\Delta_{0}}{\eta\bigl(1-L\eta\bigr)}.
\tag{9}
\]

\item[Step 7.] \textbf{Averaging.}  
Divide (9) by $T$ to obtain the claimed bound:
\[
\frac{1}{T}\sum_{k=0}^{T-1}\|G_{k}\|_{*}^{2}
\le
\frac{2\Delta_{0}}{\eta\bigl(1-L\eta\bigr)T}.
\tag{10}
\]
Since the average of non‑negative numbers dominates the minimum,
\[
\min_{0\le k<T}\|G_{k}\|_{*}^{2}
\le
\frac{2\Delta_{0}}{\eta\bigl(1-L\eta\bigr)T}
=O\!\left(\frac{1}{T}\right).
\tag{11}
\]
\end{enumerate}
Thus Theorem~\ref{thm:main} is proved. \hfill\qedsymbol

\section*{Rate Derivation (With Constants)}
Define the constant
\[
C(\eta):=\frac{2}{\eta\bigl(1-L\eta\bigr)}.
\]
For the admissible stepsize interval $(0,1/L]$, the function $C(\eta)$ is minimized at $\eta=1/L$, yielding
\[
C_{\min}= \frac{2L}{1/ L \cdot (1-1)}\;\text{(limit)}\;\to 2L.
\]
Hence a convenient choice $\eta=1/(2L)$ gives
\[
\frac{1}{T}\sum_{k=0}^{T-1}\|G_{k}\|_{*}^{2}
\le
\frac{8L\Delta_{0}}{T},
\]
i.e. an explicit $O(L\Delta_{0}/T)$ rate.

\section*{Special Cases}
\begin{itemize}
\item \textbf{Euclidean case.} $h(x)=\frac12\|x\|_{2}^{2}$, $\Breg{x}{y}=\frac12\|x-y\|_{2}^{2}$.
Then $G_{k}=\frac{1}{\eta}(x_{k}-x_{k+1})$ and $\|G_{k}\|_{*}=\|G_{k}\|_{2}$.
The theorem reduces to the standard non‑convex gradient descent guarantee.
\item \textbf{Negative entropy.} $h(x)=\sum_{i}x_{i}\log x_{i}$ on the simplex.
The dual norm becomes $\|\cdot\|_{\infty}$, and the bound holds for the
exponentiated gradient update.
\item \textbf{Strongly non‑convex setting.} If additionally $f$ satisfies the
Polyak–Łojasiewicz (PL) condition,
$f(x)-f^{*}\le\frac{1}{2\mu}\|\nabla f(x)\|_{*}^{2}$,
then (10) immediately yields a linear convergence rate
$f(x_{k})-f^{*}=O\big((1-\mu\eta)^{k}\big)$.
\end{itemize}

\end{document}